\chapter{Programmation sans verrou}\label{ch:ll:lock-free-programming}

En juillet 1997, une sonde posée à la surface de Mars se mit à redémarrer d'elle-même, perdant des données à chaque fois. Une tâche de faible priorité qui
collectait des données météorologiques détenait un verrou dont avait besoin une tâche de haute priorité gérant le bus d'information~; une tâche de priorité moyenne
s'intercalait, la tâche de haute priorité manqua son échéance, et un watchdog redémarra la machine. Le correctif, téléversé depuis
la Terre, activa l'héritage de priorité pour ce verrou (Jones, 1997). Un thread de trading qui partage un verrou avec un outil de journalisation,
un module de reporting des risques ou un rechargeur de configuration pose le même problème, sauf que l'échéance se compte en microsecondes, et que le
thread inférieur n'a pas besoin d'être interrompu par rien de plus sinistre que l'ordonnanceur du système d'exploitation. Ce chapitre traite de
du partage de données entre threads \omterm{def:ll:lock-free-programming:progress}{sans verrou}~: ce que le matériel garantit sur la mémoire, ce que C++ et Rust permettent à un programme
dire sur l'ordonnancement, la boucle \omterm{def:ll:lock-free-programming:cas}{compare-and-swap} et son piège classique, la récupération de la mémoire que d'autres threads peuvent
encore lire, et ce que «~\omterm{def:ll:lock-free-programming:progress}{sans verrou}~» promet réellement.

\section{Atomiques et modèle mémoire}

\begin{definition}[Opération atomique]\label{def:ll:lock-free-programming:atomic}
Une \emph{opération atomique}\index{opération atomique} sur un emplacement mémoire est indivisible vis-à-vis des autres threads~: aucun
thread ne l'observe à moitié faite. C++ fournit \texttt{std::atomic<T>} et Rust \texttt{AtomicU64} et ses semblables, avec
lectures, écritures, échanges, fetch-and-add et \omterm{def:ll:lock-free-programming:cas}{compare-and-swap}~; sur x86, chacun correspond à une seule instruction pour des tailles allant jusqu'à huit
octets.
\end{definition}

Une \omterm{def:ll:rust-for-low-latency:race}{course aux données} (data race, chapitre~9) est un \omterm{def:ll:cpp-for-latency-iii-the-compiler:ub}{comportement indéfini} en C++~; les atomiques sont le moyen de partager un emplacement sans en produire une. La
question qu'elles laissent ouverte est celle de l'ordre. Un processeur exécute dans le désordre (chapitre~2), garde les écritures dans un tampon avant
qu'elles atteignent le cache, et un compilateur réordonne les accès mémoire qu'il peut prouver indépendants dans un même thread~; un autre thread peut
donc voir les écritures d'un thread dans un ordre différent de celui du programme.

\begin{definition}[Ordonnancement mémoire, ordonnancement acquire--release, cohérence séquentielle]\label{def:ll:lock-free-programming:ordering}
L'\emph{ordonnancement mémoire}\index{ordonnancement mémoire} (memory ordering) d'une \omterm{def:ll:lock-free-programming:atomic}{opération atomique} précise les autres accès mémoire qu'elle ordonne.
Dans l'\emph{ordonnancement acquire--release}\index{ordonnancement acquire--release}, une lecture en acquire qui lit la valeur écrite par une
écriture en release voit toutes les écritures faites par le thread qui écrit avant cette écriture~: la paire publie les données. \emph{Cohérence
séquentielle}\index{cohérence séquentielle}, la valeur par défaut des atomiques C++, ajoute un ordre total unique sur toutes ces opérations
que tous les threads acceptent, comme si elles étaient entrelacées sur un seul processeur. Une opération relâchée est atomique mais ordonne
rien d'autre.
\end{definition}

Le schéma de publication (écrire les données, puis stocker un drapeau en release~; charger le drapeau en acquire, puis lire les données) est
le backbone de tout tampon circulaire du chapitre~12. La \omterm{def:ll:lock-free-programming:ordering}{cohérence séquentielle} est nécessaire moins souvent et coûte plus cher~; le cas classique
où l'\omterm{def:ll:lock-free-programming:ordering}{ordonnancement acquire--release} ne suffit pas est le schéma du tampon d'écritures de la \cref{fig:ll:lock-free-programming:litmus}~: chaque
thread écrit son propre drapeau et lit celui de l'autre. Sur x86, une écriture peut encore se trouver dans le tampon d'écritures du processeur quand la
lecture suivante d'un autre emplacement s'achève, si bien que les deux threads peuvent lire zéro. Seule une écriture à \omterm{def:ll:lock-free-programming:ordering}{cohérence séquentielle} (compilée
en une instruction qui vide le tampon) l'interdit.

\begin{omfigure}
\begin{tikzpicture}
\begin{axis}[width=9cm,height=5cm,ybar,bar width=16pt,ymin=0,enlarge y limits={upper,value=0.2},
  symbolic x coords={relaxed,release-acquire,seq_cst},xtick=data,xticklabels={relaxed,{release--acquire},seq\_cst},
  enlarge x limits=0.3,ylabel={$r_1 = r_2 = 0$ par million},tick label style={font=\small},label style={font=\small},
  grid=major,grid style={gray!25},nodes near coords,nodes near coords style={font=\footnotesize},table/col sep=comma]
\addplot[fill=omThm!45,draw=omThm] table[x=ordering,y=both_zero]{figdata/low-latency/11-lock-free-programming/measured_litmus.csv};
\end{axis}
\end{tikzpicture}

\omcaption{Le test litmus du tampon d'écritures~: le thread A exécute \texttt{x = 1; r1 = y}, le thread B exécute \texttt{y = 1; r2 = x},
un million de fois chacun. Sous une \omterm{def:ll:lock-free-programming:ordering}{cohérence séquentielle}, au moins une lecture voit 1~; avec des ordonnancements plus faibles, les écritures encore dans le tampon de chaque
cœur permettent aux deux de lire 0, rarement mais réellement. Mesuré sur un ordinateur portable (Intel Core Ultra 7 155H) sous WSL2, sans
cœurs isolés. Données~: \texttt{bench\_lockfree.py}.}\label{fig:ll:lock-free-programming:litmus}
\end{omfigure}

\section{Boucles compare-and-swap et problème ABA}

\begin{definition}[Compare-and-swap]\label{def:ll:lock-free-programming:cas}
Le \emph{compare-and-swap}\index{compare-and-swap} (CAS) remplace atomiquement la valeur d'un emplacement par une nouvelle valeur si et
seulement si elle est toujours égale à une valeur attendue, et indique si elle l'a fait. Une boucle CAS lit la valeur courante, calcule une nouvelle
valeur, puis relance le CAS jusqu'à ce qu'aucun autre thread n'ait modifié l'emplacement entre-temps.
\end{definition}

Toute structure \omterm{def:ll:lock-free-programming:progress}{sans verrou} est construite à partir de boucles CAS. Une pile dépile en lisant le nœud de tête et son successeur, puis fait
basculer la tête du nœud vers son successeur par un CAS~; le CAS échoue, et la boucle recommence, si un autre thread a déplacé la tête
entre-temps. L'échec qu'elle ne peut pas voir est le changement qui a été annulé.

\begin{definition}[Problème ABA]\label{def:ll:lock-free-programming:aba}
Le \emph{problème ABA}\index{problème ABA} est l'échec d'un \omterm{def:ll:lock-free-programming:cas}{compare-and-swap} à voir qu'un emplacement est passé de
$A$ à $B$, puis de nouveau à $A$, entre la lecture et le CAS~: la comparaison réussit alors que l'état à partir duquel elle a été calculée
(ici, le successeur de $A$) est périmé.
\end{definition}

\begin{omfigure}
\begin{tikzpicture}[font=\footnotesize,
  nd/.style={draw=omDef,fill=omDef!10,circle,minimum size=0.6cm,inner sep=0pt},
  fr/.style={draw=gray,fill=gray!15,circle,minimum size=0.6cm,inner sep=0pt}]
\node[anchor=south,font=\footnotesize\bfseries,align=center] at (0.9,1.9) {1. thread 1 lit\\ tête et suivant};
\node[nd] (a1) at (0,1.4) {A}; \node[nd] (b1) at (0.9,1.4) {B}; \node[nd] (c1) at (1.8,1.4) {C};
\draw[-Stealth] (a1) -- (b1); \draw[-Stealth] (b1) -- (c1);
\node[anchor=north,align=center] at (0.9,0.9) {tête = A,\\ suivant = B};
\node[anchor=south,font=\footnotesize\bfseries,align=center] at (5.3,1.9) {2. thread 2~: dépile A,\\ dépile B, empile A};
\node[nd] (a2) at (4.5,1.4) {A}; \node[nd] (c2) at (5.4,1.4) {C}; \node[fr] (b2) at (6.6,1.4) {B};
\draw[-Stealth] (a2) -- (c2);
\node[anchor=north,align=center] at (5.5,0.9) {tête = A à nouveau~;\\ B de retour dans la réserve};
\node[anchor=south,font=\footnotesize\bfseries,align=center] at (10.0,1.9) {3. thread 1~:\\ CAS(tête, A $\to$ B)};
\node[fr,draw=omThm,fill=omThm!15] (b3) at (9.5,1.4) {B}; \node[nd] (c3) at (10.6,1.4) {C};
\draw[-Stealth,omThm] (b3) -- (c3);
\node[anchor=north,align=center,omThm] at (10.05,0.9) {réussit~: la tête est\\ un nœud de la réserve libre};
\end{tikzpicture}

\omcaption{Le \omterm{def:ll:lock-free-programming:aba}{problème ABA} sur une pile \omterm{def:ll:lock-free-programming:progress}{sans verrou}. La tête vaut A les deux fois où le thread 1 regarde, si bien que son \omterm{def:ll:lock-free-programming:cas}{compare-and-swap} réussit
et installe B, que le thread 2 a déjà retiré. Un marqueur incrémenté à chaque mise à jour rend inégale la seconde comparaison avec A.
Données~: \texttt{ll\_aba.aba}, un rejeu déterministe.}\label{fig:ll:lock-free-programming:aba}
\end{omfigure}

Les remèdes classiques changent ce que l'on compare. Un marqueur (un compteur de version) stocké à côté du pointeur et incrémenté à chaque
mise à jour fait comparer les trois états $A$, $B$, $A$ comme $(A,1)$, $(B,2)$, $(A,3)$. Sur x86-64, un pointeur et un compteur tiennent dans un
\omterm{def:ll:lock-free-programming:cas}{compare-and-swap} de 16 octets~; la pile du tutoriel évite cette instruction large en rangeant les nœuds dans un tableau de taille fixe et en empaquetant un
index de 32 bits avec un marqueur de 32 bits dans un seul mot de 64 bits. Ses nœuds ne sont jamais libérés, ce qui supprime l'autre moitié du
problème.

\section{Récupération sûre de la mémoire}

\begin{definition}[Pointeur de danger, récupération par époques]\label{def:ll:lock-free-programming:reclaim}
Dans une structure \omterm{def:ll:lock-free-programming:progress}{sans verrou} dont les nœuds sont libérés, un nœud retiré par un thread peut encore être lu par un autre qui a chargé son
adresse avant ce retrait. Un \emph{pointeur de danger}\index{pointeur de danger} (hazard pointer) est un pointeur publié par chaque thread vers un nœud que le
thread s'apprête à lire~; un nœud n'est libéré que lorsqu'aucun pointeur de danger ne le désigne (Michael, 2004). Dans
\emph{récupération par époques}\index{récupération par époques} (epoch-based reclamation), les threads annoncent leur entrée dans une époque globale et leur sortie de celle-ci~; un nœud
retiré à l'époque $e$ est libéré dès que chaque thread a quitté $e$ (Fraser, 2004).
\end{definition}

Les deux reportent la libération jusqu'à ce qu'aucun lecteur ne puisse tenir le nœud, et toutes deux pèsent sur le chemin de lecture~: les \omterm{def:ll:lock-free-programming:reclaim}{pointeurs de danger} coûtent une écriture et une
barrière mémoire par nœud lu, les époques une écriture à l'entrée et à la sortie de chaque opération. Sur un chemin critique (hot path), la récupération la plus simple consiste à n'en faire aucune~:
des pools à taille fixe dont les emplacements sont recyclés mais jamais rendus au système (chapitre~6), des tampons circulaires bornés dont les cellules sont réutilisées
selon leur position (chapitre~12), et les indices étiquetés du tutoriel. La récupération compte pour les structures non bornées partagées entre
de nombreux threads, et un chemin critique ne devrait pas en avoir.

\section{Garanties de progression}

\begin{definition}[Sans verrou, sans attente]\label{def:ll:lock-free-programming:progress}
Un algorithme est \emph{sans verrou}\index{sans verrou} (lock-free) si, quel que soit l'ordonnancement des threads, un thread achève son
opération en un nombre fini de ses propres pas~: un thread qui s'arrête ne peut pas bloquer les autres. Il est \emph{sans attente}\index{sans attente} (wait-free)
si chaque thread achève son opération en un nombre borné de ses propres pas, quoi que fassent les autres (Herlihy, 1991).
\end{definition}

\begin{definition}[Inversion de priorité]\label{def:ll:lock-free-programming:inversion}
L'\emph{inversion de priorité}\index{inversion de priorité} (priority inversion) survient lorsqu'un thread de haute priorité attend une ressource détenue par un
thread de basse priorité qui ne peut pas s'exécuter, parce que des threads de priorité intermédiaire monopolisent son processeur~; le thread de haute priorité est
de fait rétrogradé sous les threads de priorité intermédiaire.
\end{definition}

Ces garanties portent sur les blocages, non sur la vitesse. Les threads d'un algorithme \omterm{def:ll:lock-free-programming:progress}{sans verrou} peuvent encore entrer en collision et réessayer, et en cas de
contention, une boucle CAS peut être plus lente qu'un verrou (\cref{fig:ll:lock-free-programming:counters})~; ce qu'elle supprime est la
possibilité qu'un thread préempté détenant un verrou fige tout le monde et, avec elle, l'\omterm{def:ll:lock-free-programming:inversion}{inversion de priorité}. Beaucoup de structures utiles
se situent entre ces catégories~: la file multi-producteurs bornée de la construction utilise une seule opération CAS par appel, \omterm{def:ll:lock-free-programming:progress}{sans
verrou}, mais son auteur signale qu'elle n'est pas \omterm{def:ll:lock-free-programming:progress}{sans verrou} au sens officiel, car un producteur arrêté entre la réservation d'une
cellule et son remplissage bloque les consommateurs de cette cellule. Les structures \omterm{def:ll:lock-free-programming:progress}{sans attente} sont plus rares et souvent plus lentes en moyenne~; un
tampon circulaire à un seul producteur et un seul consommateur (chapitre~12) est \omterm{def:ll:lock-free-programming:progress}{sans attente} sans surcoût, et c'est la structure la plus utilisée sur les chemins critiques du trading.

\begin{omfigure}
\begin{tikzpicture}
\begin{axis}[width=10cm,height=5.8cm,ymode=log,xmin=0.8,xmax=4.2,ymin=2,ymax=400,xtick={1,2,3,4},
  ytick={2,5,10,20,50,100,200},yticklabels={2,5,10,20,50,100,200},
  xlabel={threads effectuant un incrément},ylabel={ns par incrément, par thread (log)},tick label style={font=\small},
  label style={font=\small},grid=major,grid style={gray!25},
  legend columns=4,legend style={font=\footnotesize,at={(0.5,-0.25)},anchor=north,draw=none,fill=none,
  /tikz/every even column/.append style={column sep=6pt}},table/col sep=comma]
\addplot[omThm,very thick,mark=*] table[x=threads,y=mutex]{figdata/low-latency/11-lock-free-programming/measured_counters.csv};
\addplot[omProp,thick,mark=square*] table[x=threads,y=CAS_loop]{figdata/low-latency/11-lock-free-programming/measured_counters.csv};
\addplot[omDef,thick,mark=triangle*] table[x=threads,y=fetch_add]{figdata/low-latency/11-lock-free-programming/measured_counters.csv};
\addplot[omMeth,very thick,mark=diamond*] table[x=threads,y=per_thread]{figdata/low-latency/11-lock-free-programming/measured_counters.csv};
\legend{mutex,boucle CAS,fetch\_add,compteurs par thread}
\end{axis}
\end{tikzpicture}

\omcaption{Un à quatre threads comptant chacun un million d'événements dans un compteur partagé protégé par un mutex, mis à jour par une
boucle CAS ou par \texttt{fetch\_add}, ou dans des compteurs par thread alignés, additionnés à la fin. Seul le dispositif qui ne partage rien
passe à l'échelle. Mesuré sur un ordinateur portable (Intel Core Ultra 7 155H) sous WSL2, sans cœurs isolés, threads non épinglés. Données~:
\texttt{bench\_lockfree.py}.}\label{fig:ll:lock-free-programming:counters}
\end{omfigure}

Les compteurs rendent concrète la seconde leçon du chapitre. Avec un seul thread, chaque dispositif coûte quelques nanosecondes. Avec deux threads ou
plus, tout dispositif qui écrit une même ligne partagée paie à chaque mise à jour le va-et-vient de la ligne entre les cœurs (chapitre~3), une boucle CAS aussi pour ses reprises,
un mutex aussi pour la ligne propre du verrou~; les compteurs par thread, chacun sur sa propre \omterm{def:ll:memory-hierarchy-and-caches:line}{ligne de cache}, coûtent
ce qu'ils coûtent seuls. La synchronisation la plus rapide est celle dont on n'a pas besoin~: donner à chaque thread ses propres données et
les combiner à la fin, ou ne transmettre les données que dans un seul sens.

\section{Tutoriel~: une pile, une file et un test litmus}\label{tut:ll:lock-free-programming:tut}

\begin{tutorial}
\textbf{Objectif.} Construire une pile \omterm{def:ll:lock-free-programming:progress}{sans verrou} insensible au \omterm{def:ll:lock-free-programming:aba}{problème ABA}, mesurer le coût du partage d'un compteur, observer le
réordonnancement mémoire, construire une file bornée multi-producteurs multi-consommateurs. \textbf{État final~:}
les \cref{fig:ll:lock-free-programming:litmus,fig:ll:lock-free-programming:counters} et une file soumise à des tests de résistance dans deux langages.
\begin{tutsteps}
\item \textbf{Une pile étiquetée.} La tête regroupe une étiquette avec un indice de nœud~; push et pop sont des boucles CAS qui incrémentent l'étiquette.
\omcode{code/low-latency/11-lock-free-programming/cpp/ll_lockfree.hpp}{13}{43}{Une pile de Treiber sur un pool à taille fixe, à tête étiquetée.}\label{lst:ll:lock-free-programming:stack}
\item \textbf{Rejouer ABA.} \texttt{ll\_aba.aba(tagged=False)} rejoue l'entrelacement de la \cref{fig:ll:lock-free-programming:aba}
et renvoie une pile dont la tête est un nœud libéré~; avec \texttt{tagged=True}, l'opération CAS échoue.
\item \textbf{La file.} Un producteur réserve une position par une opération CAS, écrit sa cellule et la publie par une
écriture release du numéro de séquence de la cellule~; le consommateur de cette position effectue une lecture acquire.
\omcode{code/firm/mpmcq/cpp/firm_mpmcq.hpp}{30}{48}{Le côté producteur de la file MPMC bornée.}\label{lst:ll:lock-free-programming:push}
\item \textbf{Tests de résistance et sanitisation.} Les tests font tourner deux producteurs et deux consommateurs sur 300\,000 éléments et vérifient
que chaque élément arrive exactement une fois, en C++ et en Rust~; le test C++ s'exécute aussi sous le sanitiseur de threads.
\item \textbf{Mesurer} avec \texttt{python bench\_lockfree.py}.
\end{tutsteps}
\textbf{Pour aller plus loin.} Retirer l'étiquette de la tête de la pile et lancer le test d'alternance à deux threads jusqu'à ce qu'il échoue~; passer
l'écriture release de la file à relaxed et lancer le test de résistance sous le sanitiseur de threads.
\end{tutorial}

\section{Construction~: la file bornée MPMC}\label{bld:ll:lock-free-programming:mpmcq}

\begin{build}
\bfield{Rôle} La file destinée aux rares endroits où plusieurs threads produisent dans un même flux ou en consomment~: ordres
de requêtes de plusieurs stratégies vers une même passerelle, travail pour un pool de calculateurs de risque. Les chemins à un seul producteur utilisent
plutôt les tampons circulaires du chapitre~12.
\bfield{Interface} C++20~: \texttt{firm::mpmcq::Queue<T>(capacity)} avec \texttt{try\_push(v)} et \texttt{try\_pop()} renvoyant
\texttt{std::optional<T>}~; Rust \texttt{firm\_mpmcq::Queue<T: Copy>} avec \texttt{try\_push} et \texttt{try\_pop}.
\bfield{Règles} Capacité égale à une puissance de deux~; une opération CAS par appel~; aucune allocation après la construction~; l'état plein et vide
est signalé, jamais attendu à l'intérieur de la file~; positions du producteur et du consommateur sur des \omterm{def:ll:memory-hierarchy-and-caches:line}{lignes de cache} distinctes.
\bfield{Tests d'acceptation} \texttt{code/firm/mpmcq/}~: premier entré, premier sorti au fil d'un cycle plein puis vide~; deux producteurs et deux
consommateurs livrent 300\,000 éléments exactement une fois, en C++20 et en Rust, en quelques secondes.
\bfield{Pour aller plus loin} Opérations par lots qui réservent plusieurs positions par une seule opération CAS~; un enveloppeur bloquant pour les threads qui
peuvent se mettre en sommeil.
\end{build}

\begin{omsources}
\item M.~Jones, ``What really happened on Mars Rover Pathfinder'', 1997.
\item H.-J.~Boehm et S.~Adve, ``Foundations of the C++ concurrency memory model'', PLDI 2008.
\item M.~Herlihy, ``Wait-free synchronization'', \emph{ACM TOPLAS} 13(1), 1991.
\item M.~Michael, ``Hazard pointers: safe memory reclamation for lock-free objects'', \emph{IEEE TPDS} 15(6), 2004;
K.~Fraser, \emph{Practical lock-freedom}, Cambridge, rapport technique 579, 2004.
\item D.~Vyukov, ``Bounded MPMC queue'', 1024cores.
\end{omsources}

\section{Exercices}

\begin{exercise}[$\star$]\label{exo:ll:lock-free-programming:1}
À partir de la \cref{fig:ll:lock-free-programming:counters}, calculer le coût d'un \texttt{fetch\_add} partagé par incrément avec un seul thread
et avec quatre. Combien d'incréments par seconde quatre threads réalisent-ils ensemble, en compteur partagé et en compteurs par thread~?
\end{exercise}

\begin{exercise}[$\star$]\label{exo:ll:lock-free-programming:2}
Un rédacteur stocke un prix, puis lève un drapeau avec release~; un lecteur charge le drapeau avec acquire et, s'il est levé, lit le prix.
Le lecteur peut-il voir le drapeau levé et un prix ancien~? Et si l'écriture du drapeau est relaxed~?
\end{exercise}

\begin{exercise}[$\star$]\label{exo:ll:lock-free-programming:3}
Une étiquette de 32 bits est incrémentée à chaque mise à jour de la tête de la pile, à raison de 100 millions de mises à jour par seconde. Au bout de combien de temps elle revient à zéro,
et pourquoi ce rebouclage est-il (presque) sans danger~?
\end{exercise}

\begin{exercise}[$\star\star$]\label{exo:ll:lock-free-programming:4}
Dans le test du tampon de stockage, quels résultats $(r_1, r_2)$ sont possibles sous \omterm{def:ll:lock-free-programming:ordering}{cohérence séquentielle}, et quel résultat supplémentaire x86
autorise-t-il avec release--acquire~? Pourquoi la figure le montre-t-elle si rarement~?
\end{exercise}

\begin{exercise}[$\star\star$]\label{exo:ll:lock-free-programming:5}
Pourquoi la file MPMC bornée n'est-elle pas \omterm{def:ll:lock-free-programming:progress}{sans verrou} au sens strict~? Décrire le blocage.
\end{exercise}

\begin{exercise}[$\star\star$]\label{exo:ll:lock-free-programming:6}
Un thread de journalisation détient un mutex pendant \qty{2}{\micro\second} par message, 100\,000 fois par seconde, et le thread de trading prend le
même mutex une fois par ordre. Quelle fraction des ordres trouve le mutex occupé, et combien de temps peut-on attendre si le
thread de journalisation est préempté pendant un quantum d'ordonnanceur de \qty{4}{\milli\second}~?
\end{exercise}

\begin{exercise}[$\star\star\star$]\label{exo:ll:lock-free-programming:7}
\emph{Programmation.} Exécuter la pile sans étiquette de \texttt{ll\_lockfree.hpp} (mode \texttt{aba} du banc d'essai) avec deux threads
qui dépilent et repilent un million de fois chacun, vingt fois de suite. Un nœud est-il jamais perdu ou dupliqué~? Que dit le résultat
sur la validation du code \omterm{def:ll:lock-free-programming:progress}{sans verrou} par l'exécution~?
\end{exercise}

\begin{exercise}[$\star\star\star$]\label{exo:ll:lock-free-programming:8}
\emph{Trouver la faille.} «~Notre instantané du carnet d'ordres est publié avec un drapeau relaxed~; nous l'avons testé pendant une semaine sur x86 sans
jamais voir d'instantané déchiré, donc il est correct.~»
\end{exercise}

\section{Problème~: le détenteur de verrou désordonnancé}

\begin{problem}[{Problème du week-end --- un verrou partagé avec le journal}]\label{pb:ll:lock-free-programming:1}
Un thread de stratégie envoie environ 20\,000 ordres par seconde et, pour chacun, prend un mutex afin d'ajouter l'ordre à un tampon de journal que
vide un thread de journalisation. Le thread de journalisation détient le même mutex pendant \qty{3}{\micro\second} à chaque vidage, 50\,000 fois par seconde.
De temps à autre, le système d'exploitation préempte le thread de journalisation pendant qu'il détient le mutex, pour un quantum d'environ
\qty{3}{\milli\second}~; supposons que cela se produise sur 0.01~\% de ses vidages.

\medskip\noindent\textbf{Partie I --- Contention.}
\begin{enumerate}
\item Quelle fraction du temps le thread de journalisation détient-il le mutex~?
\item Quelle fraction des ordres trouve le mutex occupé, et combien de temps attendent-ils en moyenne~?
\item Quel est le temps de blocage moyen attendu par seconde pour la stratégie, en ignorant la préemption~?
\item Que dit la \cref{fig:ll:lock-free-programming:counters} du coût propre du mutex lorsque deux threads le prennent souvent~?
\end{enumerate}

\medskip\noindent\textbf{Partie II --- Préemption.}
\begin{enumerate}[resume]
\item Combien de préemptions du détenteur du verrou se produisent par seconde~?
\item Quelle fraction du temps le mutex est-il détenu par un thread de journalisation préempté~?
\item Combien d'ordres par seconde rencontrent un tel détenteur préempté, et combien de temps chacun attend-il au plus~?
\item Additionner les parties I et II~: le temps de blocage attendu par seconde.
\end{enumerate}

\medskip\noindent\textbf{Partie III --- La refonte.}
\begin{enumerate}[resume]
\item Remplacer le mutex par un anneau à producteur unique et consommateur unique~: que fait désormais la stratégie par ordre~?
\item Que se passe-t-il lorsque l'anneau est plein~?
\item Quelle garantie de progression l'anneau offre-t-il à la stratégie~?
\item Pourquoi l'héritage de priorité ne résout-il pas le problème pour un thread de trading~?
\end{enumerate}

\medskip\noindent\textbf{Partie IV --- Le verdict.}
\begin{enumerate}[resume]
\item Énoncer le \emph{résultat nommé}~: le temps de blocage moyen par seconde de la stratégie avec le mutex, puis avec l'anneau.
\item Que firent les ingénieurs de Mars Pathfinder, et pourquoi cela suffit-il là-bas~?
\item Où ailleurs dans un processus de trading un verrou peut-il se cacher~?
\item Qu'exige le \omterm{def:ll:lock-free-programming:aba}{problème ABA} d'une structure \omterm{def:ll:lock-free-programming:progress}{sans verrou} qui libère la mémoire~?
\item Comment les tests de la construction vérifient-ils la file, et que ne peuvent-ils pas prouver~?
\item Quand un mutex est-il le bon choix dans un système de trading~?
\item Quel \omterm{def:ll:lock-free-programming:ordering}{ordonnancement mémoire} la publication de l'anneau exige-t-elle~?
\item En une phrase~: qu'apporte l'approche \omterm{def:ll:lock-free-programming:progress}{sans verrou} à un chemin critique~?
\end{enumerate}
\end{problem}

\section{Questions d'entretien}

\begin{interviewq}[$\star$ \iqroles{développeur}]\label{iq:ll:lock-free-programming:1}
Quelle est la différence entre \texttt{memory\_order\_relaxed}, les opérations acquire--release et les \omterm{def:ll:lock-free-programming:atomic}{opérations atomiques} à \omterm{def:ll:lock-free-programming:ordering}{cohérence séquentielle}~?
\end{interviewq}

\begin{interviewq}[$\star\star$ \iqroles{développeur}]\label{iq:ll:lock-free-programming:2}
Expliquez le \omterm{def:ll:lock-free-programming:aba}{problème ABA} et deux manières de l'empêcher.
\end{interviewq}

\begin{interviewq}[$\star\star$ \iqroles{développeur}]\label{iq:ll:lock-free-programming:3}
Un algorithme \omterm{def:ll:lock-free-programming:progress}{sans verrou} est-il toujours plus rapide qu'un algorithme avec verrous~? Donnez un contre-exemple.
\end{interviewq}

\begin{interviewq}[$\star\star$ \iqroles{développeur}]\label{iq:ll:lock-free-programming:4}
Qu'est-ce que l'\omterm{def:ll:lock-free-programming:inversion}{inversion de priorité}, et pourquoi importe-t-elle dans un système de trading~?
\end{interviewq}

\begin{interviewq}[$\star\star$ \iqroles{développeur}]\label{iq:ll:lock-free-programming:5}
Comment testez-vous du code \omterm{def:ll:lock-free-programming:progress}{sans verrou}~? Que ne peut pas vous apprendre un test de résistance~?
\end{interviewq}

\begin{interviewq}[$\star\star\star$ \iqroles{développeur}]\label{iq:ll:lock-free-programming:6}
Concevez un compteur que de nombreux threads incrémentent et qu'un thread lit une fois par seconde, avec le coût le plus faible
pour les threads qui incrémentent.
\end{interviewq}
